\section{Conclusion}
\label{sec:conclusion}

In conclusion, the Visitor design pattern provides a robust, OOP-compliant architectural solution for scenarios where an object hierarchy is highly stable, yet the operations performed upon it are subject to frequent expansion. By elegantly decoupling algorithms from the underlying data structures, the pattern effectively upholds the Open/Closed Principle in terms of operational extensibility \cite{gof1994}.

However, as analyzed throughout this report, the Classic Visitor pattern and its advanced variants are not universal panaceas. Each architectural approach introduces distinct engineering trade-offs that must be meticulously evaluated prior to implementation:

\begin{itemize}
    \item The \textbf{Classic Visitor} provides standard structural decoupling and strict type safety but suffers from rigid circular dependencies and the moderate overhead of double virtual dispatching.
    \item The \textbf{Modern \texttt{std::variant} Visitor} (C++17) delivers unparalleled runtime performance and optimal cache locality by completely eliminating virtual method tables. It is highly recommended for modern, performance-critical C++ codebases, despite its compromise on the Dependency Inversion Principle \cite{cpprefvariant,cpprefvisit}.
    \item The \textbf{Acyclic Visitor} successfully resolves the dependency cycle issue, making it suitable for massive, highly decoupled legacy systems. Nevertheless, the severe runtime penalties associated with \texttt{dynamic\_cast} restrict its viability in latency-sensitive applications \cite{martin1997,cpprefrtti}.
\end{itemize}

Ultimately, selecting the appropriate Visitor variant—or deciding whether to employ the pattern at all—requires software architects to rigorously assess the specific constraints of their domain. If the core element hierarchy is volatile and prone to frequent additions, the Visitor pattern will inevitably induce architectural fragility and should be avoided. Conversely, for stable data structures demanding continuous operational growth (such as ASTs or Document Object Models), it remains an indispensable and highly effective tool within the modern software engineering arsenal \cite{gof1994}.
